\chapter{Introduction}

\section{Cryptography in the Automotive Domain}

Modern vehicles are no longer purely mechanical systems.
Today's cars contain upwards of hundreds of Electronic Control Units (ECUs)
responsible for functions ranging from engine management and brake control to
infotainment, over-the-air software updates, and vehicle-to-everything
communication.
Due to the sheer number of features, these systems are increasingly
interconnected.
All connections can be separated into two categories: internal ones, which are
mostly bus protocols such as CAN, LIN, and Automotive Ethernet, and external
ones such as V2X and telematics interfaces.
These circumstances increase the attack surface exposed to adversaries.
A successful intrusion into a safety-relevant ECU carries consequences that
extend well beyond data loss: it can compromise brake controllers, manipulate
steering actuators, or falsify sensor data, directly endangering human lives.

Cryptographic protection of ECU communication is therefore not optional.
Standards such as AUTOSAR SecOC (Secure Onboard Communication) and
ISO/SAE~21434 establish requirements for message authentication and
confidentiality across the vehicle network.
In practice, the algorithm most commonly used to fulfill these requirements is
the Advanced Encryption Standard (AES).
It has served as the global symmetric encryption standard since its NIST
standardization in 2001~\cite{FIPS197} and is deeply embedded in automotive
security stacks, ranging from hardware security modules (HSMs) to communication
middleware.

\section{Cryptographic Requirements of Modern ECUs}

ECUs span a wide spectrum of computational capability.
On the high end are domain controllers managing ADAS functions or running an
automotive-grade operating system; these may be equipped with multi-core
processors, dedicated cryptographic co-processors, and several megabytes of RAM.
At the other end are simple body control modules, sensor nodes, and actuator
controllers built around low-cost 32-bit microcontrollers with RAM measured in
tens of kilobytes and no hardware cryptographic accelerator.
On these devices the computational cost of cryptographic operations becomes a
real concern.

For such constrained ECUs, the cryptographic primitive must satisfy several
requirements simultaneously: it must provide authenticated encryption with
associated data (AEAD) to guarantee confidentiality, integrity, and authenticity
in a single pass; it must execute within a time budget that does not interfere
with real-time control tasks; and its implementation must fit within the
available Flash and RAM without displacing application code.
These requirements motivate a closer examination of whether AES-128-GCM ---
designed for general-purpose hardware --- remains the most appropriate choice
for every node in the vehicle network, or whether purpose-built lightweight
alternatives deserve consideration.

\section{The Limits of AES on Low-End Microcontrollers}

AES-128-GCM is the dominant AEAD construction in modern security protocols and
is well understood with respect to both security and implementation.
On hardware that includes a dedicated AES accelerator, it executes efficiently
and introduces negligible overhead.
However, a large portion of automotive microcontrollers --- particularly those
in cost-sensitive, high-volume applications --- are based on simple 32-bit cores
such as the ARM Cortex-M0, which provide no hardware cryptographic acceleration.
On these platforms, a software implementation of AES-128-GCM must perform the
full AES key schedule and block cipher in software, and the GCM authentication
step requires a 128-bit polynomial multiplication (GHASH) that maps poorly onto
a minimal instruction set without hardware multiply-accumulate support.
The result can be an implementation that is slower, larger in code size, or more
energy-intensive than the available budget permits.

This raises a practical engineering question: for constrained ECUs where AES
hardware acceleration is absent, is there a standardized alternative that offers
better performance characteristics without sacrificing security?

\section{The NIST Lightweight Cryptography Standardization Process}

Recognizing the growing need for cryptographic primitives suited to constrained
devices, NIST initiated a Lightweight Cryptography (LWC) standardization project
in 2018~\cite{NISTLWCCall}.
In February 2023, NIST announced \mbox{ASCON} as the winner~\cite{ASCON2023}.
ASCON is a family of lightweight cryptographic algorithms based on a sponge
construction, designed from the ground up for efficient software implementation
on small processors.
Its selection provides the embedded and automotive security community with a
standardized, well-analyzed alternative to AES-GCM for use cases where resource
constraints are the primary concern.

Despite this standardization, ASCON adoption in automotive and industrial
embedded systems remains limited.
AES continues to dominate deployed security stacks, in part due to inertia,
existing hardware support, and a lack of quantitative performance data on
representative automotive-class microcontrollers.
This work aims to contribute directly to closing that gap.

\section{Objective and Research Questions}

This paper presents a quantitative benchmark comparing Ascon-AEAD128 and
AES-128-GCM on two ARM Cortex-M microcontrollers representative of the
low-to-mid range of the automotive ECU spectrum: the NUCLEO-F072RB
(Cortex-M0, 48\,MHz) and the NUCLEO-L476RG (Cortex-M4F, 80\,MHz).
Both boards serve as a simulation of the resource conditions found in
constrained automotive ECUs.
All measurements use software-only implementations to ensure a fair,
hardware-accelerator-independent comparison applicable to the broadest range of
real-world deployments.

The following research questions guide the study:
\begin{enumerate}\tightlist
  \item \textbf{Performance:} How do Ascon-AEAD128 and AES-128-GCM compare in
        execution time and throughput across a range of message sizes on each
        platform?
  \item \textbf{Flash footprint:} What is the Flash code-size footprint of each
        implementation relative to the constraints of a typical automotive
        microcontroller?
  \item \textbf{Optimization sensitivity:} How does compiler optimization level
        affect the relative performance of the two algorithms?
  \item \textbf{Architecture dependence:} Does the performance gap between the
        two algorithms differ between the Cortex-M0 and the Cortex-M4F, and what
        does this imply for ECU selection?
\end{enumerate}

\section{Structure of This Paper}

The remainder of this paper is organized as follows.
Chapter~2 provides the technical background on AES-128-GCM, the ASCON family,
and the target hardware platforms.
Chapter~3 describes the benchmark methodology, including the measurement setup,
timing approach, library selection, and build configuration.
Chapter~4 presents the measurement results.
Chapter~5 discusses the findings in the context of the research questions and
draws practical conclusions for algorithm selection on constrained automotive
hardware.
Chapter~6 concludes the paper and outlines directions for future work.
